
\subsection{Texture Resonance for Style Alignment}
To validate the effectiveness of the proposed Texture Resonance Retrieval (TRR) against traditional Vector Similarity, we conducted a comparative experiment (Experiment 4) using synthetic audio data. We define the retrieval error as the L2 distance between the retrieved effect parameters and the ground truth parameters for a given style query.

Table \ref{tab:exp4} presents the retrieval error for two distinct musical styles: ``Stadium Rock'' (characterized by clean, spatial tones) and ``Punk Rock'' (characterized by aggressive distortion).

\begin{table}[h]
\centering
\caption{Comparison of Texture Resonance vs. Vector Baseline (Lower is Better)}
\label{tab:exp4}
\begin{tabular}{@{}lccc@{}}
\toprule
\textbf{Test Case} & \textbf{Vector Baseline} & \textbf{Texture Resonance (Ours)} & \textbf{Result} \\ \midrule
Stadium Rock Lead & \textbf{7.76} & 8.10 & Parity \\
Punk Rock Raw & 16.93 & \textbf{7.96} & \textbf{TRR Wins} \\ \midrule
\textbf{Average} & 12.35 & \textbf{8.03} & \textbf{TRR Wins} \\ \bottomrule
\end{tabular}
\end{table}

The results demonstrate the robustness of TRR. While the Vector Baseline achieved parity in the standard case, it failed significantly (Error: 16.93) in capturing the ``Punk Rock'' style, likely due to the mean-pooling of wav2vec features smoothing out the discriminative ``aggressive'' texture. In contrast, TRR maintained a low error rate (Error: $\approx$ 8.0) across diverse styles, confirming that the Gram Matrix representation effectively captures the structural ``resonance'' of audio texture, independent of basic spectral averages.
